% p.25 — الباب الثالث العشر في اعان يعين (تابع): المنفيات والمؤكدات
\mspage{25}
% Layout: six tall columns (right-to-left), each a stack of boxed blocks of unequal height.
% Every sub-heading box holds the heading AND, beneath it, the underlying (pre-iʿlāl) form,
% e.g. «مجهول بلم / لَمْ يُعْوَنْ»; the regular forms follow outside the box. Row alignment is
% approximate: all six columns run to 21 lines, so rows are aligned 1:1; the \cline marks
% reproduce the printed box boundaries inside each column.
% Column 1 continues the مضارع بلم (معروف) block begun on p.24. The running header line of the
% frame is blank (smudged) on this page.
\begin{sarf}{|c|c|c|c|c|c|}
لَمْ تُعِيْنَا & لَمْ أُعَنْ & لَنْ يُعَانَا & لَتُعِيْنَنَّ & لَتُعَنَّانِّ & \thead{أمر حاضر بالنون}\\
لَمْ تُعِيْنُوْا & لَمْ نُعَنْ & لَنْ يُعَانُوْا & لَتُعِيْنَانِّ & لَأُعَانَنَّ & أَعْوِنَنَّ\\\cline{2-2}\cline{6-6}
لَمْ تُعِيْنِي & \thead{مضارع بلن} & لَنْ تُعَانَ & لَتُعِيْنُنَّ & لَنُعَانَنَّ & أَعِيْنَنَّ\\\cline{5-5}
لَمْ تُعِيْنَا & لَنْ يُعْوِنَ & لَنْ تُعَانَا & لَتُعِيْنِنَّ & \thead{منفي بالنون} & أَعِيْنَانِّ\\\cline{2-2}
لَمْ تُعِنَّ & لَنْ يُعِيْنَ & لَنْ يُعَنَّ & لَتُعِيْنَانِّ & لَا يُعْوِنَنَّ & أَعِيْنُنَّ\\\cline{5-5}
لَمْ أُعِنْ & لَنْ يُعِيْنَا & لَنْ تُعَانَ & لَتُعِنَّانِّ & لَا يُعِيْنَنَّ & أَعِيْنِنَّ\\\cline{5-5}
لَمْ نُعِنْ & لَنْ يُعِيْنُوْا & لَنْ تُعَانَا & لَأُعِيْنَنَّ & \thead{مجهول بالنون} & أَعِيْنَانِّ\\\cline{1-1}
\thead{مجهول بلم} & لَنْ تُعِيْنَ & لَنْ تُعَانُوْا & لَنُعِيْنَنَّ & لَا يُعْوَنَنَّ & أَعِنَّانِّ\\\cline{4-4}\cline{5-5}\cline{6-6}
لَمْ يُعْوَنْ & لَنْ تُعِيْنَا & لَنْ تُعَانِي & \thead{مجهول بالنون} & لَا يُعَانَنَّ & \thead{نهي حاضر بالنون}\\\cline{1-1}\cline{5-5}
لَمْ يُعَنْ & لَنْ يُعِنَّ & لَنْ تُعَانَا & لَيُعْوَنَنَّ & \thead{أمر غائب بالنون} & لَا تُعْوِنَنَّ\\\cline{4-4}\cline{6-6}
لَمْ يُعَانَا & لَنْ تُعِيْنَ & لَنْ تُعَنَّ & لَيُعَانَنَّ & لِيُعْوِنَنَّ & لَا تُعِيْنَنَّ\\\cline{5-5}
لَمْ يُعَانُوْا & لَنْ تُعِيْنَا & لَنْ أُعَانَ & لَيُعَانَانِّ & لِيُعِيْنَنَّ & لَا تُعِيْنَانِّ\\\cline{5-5}
لَمْ تُعَنْ & لَنْ تُعِيْنُوْا & لَنْ نُعَانَ & لَيُعَانُنَّ & \thead{مجهول بالنون} & لَا تُعِيْنُنَّ\\\cline{3-3}
لَمْ تُعَانَا & لَنْ تُعِيْنِي & \thead{مضارع بالنون} & لَتُعَانَنَّ & لِيُعْوَنَنَّ & لَا تُعِيْنِنَّ\\\cline{5-5}
لَمْ يُعَنَّ & لَنْ تُعِيْنَا & لَيُعْوِنَنَّ & لَتُعَانَانِّ & لِيُعَانَنَّ & لَا تُعِيْنَانِّ\\\cline{3-3}\cline{5-5}
لَمْ تُعَنْ & لَنْ تُعِنَّ & لَيُعِيْنَنَّ & لَيُعَنَّانِّ & \thead{نهي غائب بالنون} & لَا تُعِنَّانِّ\\\cline{6-6}
لَمْ تُعَانَا & لَنْ أُعِيْنَ & لَيُعِيْنَانِّ & لَتُعَانَنَّ & لَا يُعْوِنَنَّ & \thead{تمرين}\\\cline{5-5}\cline{6-6}
لَمْ تُعَانُوْا & لَنْ نُعِيْنَ & لَيُعِيْنُنَّ & لَتُعَانَانِّ & لَا يُعِيْنَنَّ & أَجَابَ\\\cline{2-2}\cline{5-5}
لَمْ تُعَانِي & \thead{مجهول بلن} & لَتُعِيْنَنَّ & لَتُعَانُنَّ & \thead{مجهول بالنون} & أَقَامَ\\
لَمْ تُعَانَا & لَنْ يُعْوَنَ & لَتُعِيْنَانِّ & لَتُعَانِنَّ & لَا يُعْوَنَنَّ & أَنَارَ\\\cline{2-2}\cline{5-5}
لَمْ تُعَنَّ & لَنْ يُعَانَ & لَيُعِنَّانِّ & لَتُعَانَانِّ & لَا يُعَانَنَّ & أَنَابَ\\
\end{sarf}
